% page 571 of the facsimile (image p-570 of the scan)
% block 167.6 x 242.0 mm ; left margin 34.4 mm
\pgc{12.700mm}{0.000mm}{
\l{}
\l{da ulo); (b) Recevar la tota efekto de ulo sen shancelar\cel{3}563}
\l{de lo. - EFIS.}
\l{\cel{1}\sou{supozar}.\cel{2}(\sou{trans.})\cel{2}Admisar kom realigita en la tempo-punto di}
\l{\cel{4}sua dico (to quon onu imaginas esar tala). - EFIS.}
\l{}
\l{supozitorio. (medic.)\cel{3}Kono oblonga ek substanco medikamentaa}
\l{\cel{4}solida, quan onu igas penetrar aden la anuso por eventigar}
\l{\cel{4}kaki. - DEFIS.}
\l{}
\l{\cel{1}\sou{supre}. En objekto vertikala : en la parto qua distas maxime de}
\l{\cel{4}lua pedo, de lua bazo. (\sou{anke metaf.)}\cel{2}- I.}
\l{}
\l{\cel{1}\sou{supresar}. (\sou{trans}.) Desaparigar, partikulare per impedar lua ma-}
\l{\cel{4}nifesto. - EFIS.}
\l{}
\l{\cel{1}\sou{sur}.\cel{2}Prepoziciono qua indikas la situeso di kozo relate a la ko-}
\l{\cel{3}zo qua esas an lu, interkontakte, en sama direciono vertikala :}
\l{\cel{4}(a) relate to quo portas lu ; (b) relate to quon lu kovras o}
\l{\cel{4}tegas.}
\l{}
\l{\cel{1}\sou{suro}. (\sou{anat.}) En la gambo, ta parto mola quan formacas, ye la}
\l{\cel{4}extremajo, dope, la saliajo ek la muskuli jemela e la muskulo}
\l{\cel{4}solea. - L.}
\l{}
\l{surda. Di qua la oreli ne perceptas, o perceptas nur tre desfaci--}
\l{\cel{4}le, la soni, pro la desaparo, o la feblesko tre granda, di}
\l{\cel{4}la audo-senso. - EFIS.}
\l{}
\l{\cel{1}\sou{surfaco}. I. Extensajo extera di korpo. - II. (\sou{geom}.)To quo li-}
\l{\cel{4}mitizas korpo en spaco, ed havas nur du dimensioni : longeso}
\l{\cel{4}e larjeso. - EFIS.}
\l{}
\l{surjeto. (sutado).Sut-uno ube la agulo trairas du stofo-peci po--}
\l{\cel{4}zite bordo-an-bordo, igante sempre la filo pasar super la du}
\l{\cel{4}bordi unionita. - FI.}
\l{}
\l{\cel{1}\sou{surkruto}. Kaulo hachita, qua fermentacis en sal-aquo. - D.}
\l{}
\l{\cel{1}\sou{surmenar}. (trans.) Fatigegar (kavalko-bestio, charjo-bestio, e c.}
\l{\cel{4}per igar lu irar tro rapide o dum tempo tro multa. (\sou{anke}}
\l{\cel{4}\sou{metafore}). - F.}
\l{}
\l{surmuleto. (zool.) Maro-fisho. - L. Mullus surmuletus. - DEF..}
\l{}
\l{\cel{1}\sou{surogato}. To quo povas remplasar altra substanco (farmaciala,}
\l{\cel{3}tintala, e c.) pro havar ula karakterizivi analoga, ula efektii}
\l{\cel{4}analoga. - DEIR.}
\l{}
\l{\cel{1}\sou{surpliso}. (\sou{religio krist.}) Vesto ek teoc plisuroza quan +weras,}
\l{\cel{4}sur sua sutano, la sacerdoti (en la konfeseyo ed en la proce-}
\l{\cel{4}sioni funerala), la diakoni, la kantori. - EFS.}
\l{}
\l{\cel{1}\sou{surprizar.}\cel{1}(\sou{trans}.)I. Prenar (ulu, ulo) per arivar ne-expektite.}
\l{\cel{4}- II. Trublar (ulu) per arivar ne-expektite. - III. Frapar la}
\l{\cel{3}psiko per ulo ne-expektita. - EFIS..}
}
